\section{Introduction}
\label{sec:intro}

A safety evaluation that finds nothing is the most common result in AI safety, and the least
interpretable. ``We could not elicit the behaviour'' can mean the model is robust, or it can
mean the experiment was not sensitive enough to see the failure. Static protocols cannot tell
the two apart, and the difference matters: this sentence appears in model cards, audit reports,
and deployment decisions.

The literature already contains proof that the second reading is often the correct one. In the
JailbreakBench artifacts, PAIR~\cite{chao2023pair} scores $0.00$ attack success against
\llamaseven while prompt-plus-random-search scores $0.90$ on the \emph{same} 100 behaviours of
the \emph{same} model~\cite{chao2024jbb,andriushchenko2024adaptive}. A lab that ran only PAIR
would have published a false robustness claim. Similar reversals recur across the alignment
literature: jailbreaks reported as ineffective against unlearning succeed when applied more
carefully~\cite{lucki2024unlearning}, and defenses reporting single-digit attack success under
automated single-turn attacks fall to human multi-turn red teaming more than 70\% of the
time~\cite{li2024mhj}.

The natural response, and the hypothesis we test, is that protocols should be \emph{adaptive}:
treat a null as actionable feedback and use it to trigger experimental redesign. This idea is
widely endorsed and, so far, untested as such.

\para{The gap.} Adaptivity has been argued for \emph{attacks}~\cite{andriushchenko2024adaptive}
and for \emph{forecasting} rare behaviours~\cite{jones2025forecasting}, but never with the
protocol as the independent variable and total compute held fixed. Every published adaptivity
result is confounded with spending more queries, so ``adaptive beats static'' may only mean
``more compute beats less compute.'' Worse, the word ``adaptive'' silently conflates two very
different moves: adjusting parameters \emph{inside} a design family, and \emph{switching} design
family. And nobody scores the null branch, \ie what an adaptive protocol should output when it
still finds nothing. That branch is exactly the half of the claim that safety cases rest on.

\para{Our approach.} We make the protocol the independent variable and hold the budget constant
by construction. Each substrate's outcome matrix is evaluated exhaustively once, and protocols
are then \emph{replayed} against the cached matrix under an identical budget $B$. This removes
the compute confound, makes every comparison paired on the same behaviours, and lets us report
whole detection-versus-budget curves rather than single points. We run four experiments on three
substrates: 437k per-iteration random-search traces released by
\citet{andriushchenko2024adaptive}; the official JailbreakBench cost-to-discovery
matrix~\cite{chao2024jbb} of 5 attacks $\times$ 4 target models $\times$ 100 behaviours; and a
live grid on three open-weights instruct models run locally under a real anti-sycophancy
intervention, using sycophantic answer-flipping on \gsmk~\cite{cobbe2021gsm8k} as a non-harmful
failure mode.

The critical design choice is our comparator. Most work contrasts an adaptive protocol with a
single static design, which cannot separate ``adaptivity works'' from ``design diversity
works.'' We add \fixedsplit: same budget, same set of design families, but the switch points are
placed on a blind schedule that ignores the data. Our adaptive arm \ladder is built so that
disabling its trigger recovers \fixedsplit exactly, so the two arms differ \emph{only} in
whether the null is used as a trigger.

\para{What we find.} The motivating phenomenon is real and extreme: the default design detects
$0/50$ behaviours at every budget up to $10{,}000$ iterations, while a protocol spending the
identical budget across three design families reaches $1.00$ (McNemar $p<10^{-6}$, Cohen's
$h \approx 2.4$). But against \fixedsplit, the null-as-trigger mechanism gains nothing
($\Delta = +0.02$, $p = 0.73$) and loses at large budgets ($\Delta = -0.12$, 95\% CI
$[-0.22, -0.04]$). We trace this to a mechanistic fact: conditional on a run still being null at
iteration $t$, the continuous progress signal predicts eventual success at pooled AUC
$0.45$--$0.61$, with every confidence interval covering $0.5$. Adaptivity at the \emph{portfolio}
level does work, beating an equally diversified static schedule by $+0.65$ behaviours on GPT-4
and by $+0.72$ live on \qwenseven. Finally, nominally 95\% upper bounds on residual vulnerability
achieve 63--67\% empirical coverage, and an anytime-valid correction for optional stopping moves
coverage only from $0.638$ to $0.671$.

\para{Contributions.}
\begin{itemize}[leftmargin=*,itemsep=0pt,topsep=2pt]
    \item We propose a compute-matched replay methodology that makes the evaluation protocol the
    independent variable, and we instantiate it on three substrates covering replayed attack
    traces, a benchmark cost matrix, and live open-weights models.
    \item We decompose ``adaptive'' into within-search escalation and portfolio-level allocation,
    and show with a trigger-ablated comparator that the entire measured benefit comes from
    searching multiple design families, not from the null-as-trigger mechanism.
    \item We test the mechanism directly by measuring the diagnosticity (AUC) of the progress
    signal on the surviving-null population, and find it indistinguishable from chance at all
    seven checkpoints.
    \item We score the null branch and show that nominally 95\% bounds cover 63--67\% of the
    time, that anytime-valid correction does not repair this, and that on total nulls tail
    extrapolation is undefined 100\% of the time.
    \item We replicate the whole contrast live on \qwenseven, \qwensmall and \phimini, where a
    static protocol using the natural first design measures a 3.7\% flip rate and would report
    that the anti-sycophancy intervention holds, while 85\% of items flip under a different
    pressure family.
\end{itemize}

\Secref{sec:related} positions the work, \secref{sec:method} describes the substrates and
protocol arms, \secref{sec:results} reports the four experiments, and
\twosecrefs{sec:discussion}{sec:conclusion} discuss implications and limitations.
